\section{语料与数据}

\subsection{语料}本文的语料是一本 156 章、约 149 万字的通俗明代通史（2006--2009 年间分七部连载出版），
原件不随论文发布，原因与试点篇相同：版权约束。全书 156 章共约 149 万字，平均每章约 9558 字。
抽取产物是一个四层结构：1231 个人物、563 个地点、1106 件事件、2324 条关系，全部带章节归属字段。

\subsection{与本文直接相关的字段}
\begin{itemize}
  \item $\texttt{relations[].category}$：关系类别，共 12 类。本文的主变量「政治攻讦」为 193 条（8.3\%），
        对照变量「敌对冲突」321 条（13.8\%）、「君臣从属」397 条（17.1\%）。
  \item $\texttt{relations[].rel}$：关系的自然语言描述（如「弹劾」「告发」「监视」「构陷」），
        本文用它来统计信息型动作词。
  \item $\texttt{relations[].chapter}$：该关系来自哪一章。此处有一个必须先说清的口径问题——
        2324 条关系里只有 1385 条（59.6\%）带真实章节号，其余 939 条（40.4\%）标记为「推导」，
        即由事件共现自动补出、没有章节出处。
  \item $\texttt{characters[].role / faction}$：人物角色描述与派系标签。厂卫系、内廷宦官两组
        就是按 $\texttt{role}$ 与 $\texttt{faction}$ 中是否出现指定机构名来辨认的（见方法节的机构辨认）。
\end{itemize}

\subsection{两条口径铁律（违反会静默出错）}
第一，$\texttt{relations[].from/to}$ 存的是**裸人物名**（如「朱棣」），不是「person:朱棣」形式，
按名字查度数，做前缀映射会把 Spearman $\Rho$ 直接变成 0.00。
第二，凡做截断桶（`min(v,N)`）必须改名或另立 `geN` 汇总桶；沿用被截断的键名曾造出「度数 = 10 的有 143 人」
的假象（真值 7 人）。本文所有分箱都只报真实区间。

\subsection{书内证据的核验纪律}正文引用的每处书内原句都经原书全文逐字命中：
「它却有一个诡异的名字」1 次、「短短十几天内朝廷人人皆知」1 次、「我们差点忘记了那个告密的涂节」1 次、
「卢忠是锦衣卫」1 次、「王世贞被列入了严嵩的黑名单」1 次、「宋忠派人传播谣言」1 次、
「谣言传得满天飞」1 次、「就没有互联网和手机短信的当年」类表述按「在没有互联网和手机短信的当年」命中 1 次。
候选词「内厂」「厂卫」「揭帖」「榜文」「侦缉」「暗访」「风闻」「耳语」在原书中**命中数为 0**，
全部弃用——其中「厂卫」是最容易想当然的一个组合，书中从不并称。

\subsection{关系类别体系是内造的，不是史实分类}12 类关系与 11 类事件都是抽取时人工设定的枚举，
不是明代官制或史书的原生分类。因此下文所有比例都只能表述为「本书语料在本抽取口径下」，
不能表述为「明代某现象的比例」。这一条同时构成本文限制节的一条。
